\begin{auxlemma}[Separating a scalar by Weierstrass division]\label{lem:short-arc-content}
Let \(O\) be a coefficient DVR with fraction field \(K\),
\(v_p(p)=1\), and let \(J\subset\mathcal R=O[[q]]\) be generated by
a distinguished polynomial. Take \(D\in\mathcal R\) and
\(f\in\operatorname{Frac}(\mathcal R)\). Assuming that the residue
\(\bar D\in\mathcal R/J\) is nonzero, denote by \(c_J(D)\) the minimum of the
valuations of its Weierstrass remainder coefficients. Let \(k,d\)
be nonnegative integers with \(c_J(D)\le d\). If
\(Df\in\mathcal R\), \(x\in K\), and \(p^k Df-Dx\in J\), then there is an integral constant satisfying
\[
 \exists x'\in O:\qquad D(p^d f-x')\in J.
\]
\end{auxlemma}
\begin{proof}
Write \(H=\mathcal R/J\), a free \(O\)-module with the monomials
below \(\deg J\) as basis, where \(\deg J\) is the degree of the
distinguished generator. In \(H\otimes_O K\) the given equality
implies
\[
 x\bar D=p^k\overline{Df}\in p^k H,\qquad
 v_p(x)+c_J(D)\ge k .
\]
Thus \(x'=p^{d-k}x\) is integral, including when \(k>d\).
Both \(p^d Df\) and \(x'D\) lie in \(\mathcal R\), and their
images agree after tensoring \(H\) with \(K\). Freeness makes
\(H\to H\otimes_O K\) injective, which proves the assertion.
\end{proof}

\begin{lemma}[Two-letter lifting on a short arc]\label{lem:two-letter-wach-lift}
In the models of Proposition~\ref{prop:short-arc-coordinates},
write \(e_2\) for the exponent of a two-letter correction. One can achieve
\[
 e_2\le15L+1+\eta B,
\]
where \(L\) is the local size from the short-arc comparison and \(B\) its
exceptional indicator.
\end{lemma}

\begin{proof}
\label{par:short-two-levels}
Use \(u=s_{c-1}\) for the small size and \(v=s_c\) for the large
one. Put \(p=\ell,\ L=S_g\) within this calculation. Denote by
\(A_i\) the ceiling of the least ordinary Wach image exponent in
\(\pi\) at \(i\), divided by \(p-1\) before taking the ceiling.
Then \(X^{A_i}\) kills the tame-invariant cokernel. In signed
boxes of widths \(s_i+S+\eta\), the recurrence of
Paragraph~\ref{par:wach-cut-colength} gives
\begin{equation}\label{eq:short-A}
 A_i\le(s_{i+1}+S+\eta)/p+\tfrac52 L.
\end{equation}
Here one can replace \(5/2\) by \(2\) if \(i+2\) is small.

\proofstep{Compute the denominators of the two neighboring level evaluations.}
To check the constants, write \(d=p-1,\ H=h+\eta\),
\(W_j=s_j+S+\eta\) and \(c_H=\lceil1+\log_p(1+H)\rceil\).
The ordinary image exponents satisfy
\(s_i^{\rm im}\le W_{i+1}+c_H+3+s_{i+1}^{\rm im}/d\),
and are at most \(H\) by the signed twist. Thus
\[
 A_i-W_{i+1}/p\le 1+(c_H+3)/d+H/(pd)+H/d^2.
\]
The general bound uses \(H=h+\eta\), \(H+1\le2h\), and
\(p/d\le1.1\).

If \(i+2\) is small, iterating the recurrence
once replaces its last summand by
\[
 \frac{H}{d^2}\quad\rightsquigarrow\quad
 \frac{w+S+\eta+c_H+3}{d^2}+\frac{H}{d^3}.
\]
The respective coefficients of \(h/p^2\) are bounded by
\[
 \frac{p}{d}+\frac{p^2}{d^2}\le2.31,\qquad
 \frac{p}{d}+\frac{p^2}{d^3}\le1.24 .
\]
The \(\eta/p^2\) terms are absorbed in \((w+S+\eta)/p\).
Together with \(c_H\le2+\log_p(2h)\), these give the factors
\(5/2\) and \(2\) in the two versions of \eqref{eq:short-A}.
\label{cor:inert-short-two-1}

Apply Lemma~\ref{lem:short-arc-content} with
\(P=P_1=p\varepsilon+Uq^p\), where \(\varepsilon,U\in\mathcal R^\times\).
We estimate the contents for the two possible denominators.
For \(D=q^v,J=(P^y)\), put \(t=\lfloor v/p\rfloor+1\).
Multiply by \(q^{pt-v}\) and expand
\[
 (q^p)^t=U^{-t}(P-p\varepsilon)^t.
\]
Modulo \(P^y\) the
content of this power is the minimum over \(j<y,\ j\le t\)
of \(t-j+v_p\binom tj\). After division by the lowest
uniformizer power, the least \(j\) attaining it has strictly
smaller order in the residue parameter than the other surviving
terms and than \(P^y\). This proves the minimum formula even for
nonconstant units.

Since multiplication by \(q^{pt-v}\) can only
increase content in the first quotient, any admissible index gives
an upper bound for \(c_{(P^y)}(q^v)\).
For \(D=P^u,J=(q^C)\), use the binomial expansion of
\((p\varepsilon+Uq^p)^u\) modulo \(q^C\). The two bounds are
\begin{equation}\label{eq:two-letter-wach-contents}
 \begin{aligned}
 c_{(P^y)}(q^v)&\le t-i_1+v_p\binom{t}{i_1},
       & i_1&=\min(t,y-1),\\
 c_{(q^C)}(P^u)&\le u-i_2+v_p\binom{u}{i_2},
       & i_2&=\min(u,\lfloor(C-1)/p\rfloor).
 \end{aligned}
\end{equation}
For \(pu<C\) the second line has \(i_2=u\) and the content is zero.
For the binomial terms in \eqref{eq:two-letter-wach-contents}, the carry formula gives \(v_p\binom ti\le\log_p t\).
In particular the two bounds are at most
\[
 \max(0,v/p-y+2+\log_p v)+1,\qquad
 \lceil\max(0,u-C/p+2+\log_p u)\rceil ,
\]
respectively. In the first bound the case \(i=t\) contributes
zero. All contents are normalized by \(v_p(p)=1\).

Suppose first that the large terminal letter is positive. If
\(p(u+S+\eta)+\eta<v\), construct a filtered source constant.
On the target take the coordinate \(P^{u+S}/q^v\) in direction
\(d\), realize it at cost \(\sigma\) and lift at the cut at
cost \(R\). Express the lift in source horizontal coordinates
at the cost of a further \(\sigma\). Multiply these coordinates
by \(q^{v+S}P^\eta\), obtaining integral series. In fact every
nonzero filtered summand has exponents at least \(-(v+S)\)
and \(-\eta\) at the two levels. Divisibility after rational
expansion descends by saturation. In the coordinate \(q\) take
the coefficient of degree
\[
 j=S+p(u+S+\eta).
\]

A filtered summand with level-zero exponent \(\alpha\) can
contribute only if
\[
 v+S+\alpha\le j.
\]
This forces
\(-\alpha\ge\eta+1\); the gap puts its direction in
\(\mathcal F_0^v\). The constant matrix of the map sends this
coefficient to the same coefficient of
\(p^{R+2\sigma}q^S P^{u+S+\eta}d\). Indeed the coordinate error
after clearing these denominators is divisible by \(q^h\) and
\(j<h\). The coefficient of degree \(p(u+S+\eta)\) of
\(P^{u+S+\eta}\) is a unit, by reduction modulo the maximal
ideal. Dividing by this unit supplies the filtered preimage of \(p^{R+2\sigma}d\). The additional coordinate comparison in the second rule raises the exponent to
\[
 (R+2\sigma)+\sigma=R+3\sigma.
\]


In the complementary range take \(D=q^v\), \(y=u+S\). The
first content bound separates the constant at the negative level
from the remainder. Each is lifted at cost \(R+\sigma\) by the
two rules. Writing the extra cost as \(c_{\rm cont}\), the two
ranges give
\[
 c_{\rm cont}\le
 \begin{cases}
    3+\gamma_h,& v<py,\\
    \eta+\eta/p+3+\gamma_h,& py\le v\le p(y+\eta)+\eta .
 \end{cases}
\]
In the second range, combining the two bounds on \(v\) with
\(b_0\le v\le b_0+\eta\) gives
\[
 b_0/p-S-\eta-\eta/p\le u\le(b_0+\eta)/p-S,
\]
the prescribed band.

Now suppose the first small letter is positive. Put
\(A=A_{c-1}\). When \(u\ge A+\eta+1\), take the partial
arc of length one at this letter and the coordinate \(q^{-u}d\)
there, with \(d\) read before the last Frobenius transport.
Realize it at cost \(\sigma\), lift \(X^A\) times this integral
vector, and read the source coordinates at cost \(\sigma\).
Multiply by \(q^G\) and take the coefficient in degree \(G-u+A\).

The degree \(G-u+A\) is below \(G-\eta\), so the degree and gap test places the coefficient in \(\mathcal F^u\). The target coefficient in the chosen degree is one, as
\[
 [q^{A+G-u}](q-p)^A q^{G-u}=1.
\]
The precision includes \(q^{2G}\) after denominators are cleared, so the image of the coefficient is exactly \(p^{2\sigma}d\). Tame averaging costs a unit.
Transport forward and apply the second rule at pole set
\(\{1\}\). The cost is \(3\sigma\).

When \(u\le A+\eta\), apply the second content bound with
\(D=P^u,\ C=v+S\), and lift the separated remainder and
constant at cost \(R+\sigma\). There is no extra cost for
\(pu<C\). Otherwise \eqref{eq:short-A} with its improved
constant gives
\[
 u-C/p\le\eta+\eta/p+2L.
\]

If the band is absent, \(pu\ge C\ge b_0\) and
\(pu-b_0>2p\eta\). But
\[
 u-b_0/p\le \eta+(2\eta+S)/p+2L,
\]
so \(\eta<4L\) using
\((2\eta+S)/p\le L\) (since \(w>\eta\)).
Thus the two-letter exponent in all cases is at most
\(15L+1+\eta B\). \label{cor:inert-short-two-2}
\end{proof}

